\section{Introduction}\label{sec:introduction}

An isotropic ferromagnet has no preferred direction in zero field.
In a finite periodic volume, its Gibbs state is therefore rotation
invariant even at temperatures where magnetized equilibrium states exist.
The natural description of the ordered phase is a magnetization of fixed
length and uniformly distributed direction. Rotation invariance alone
does not establish this description: it allows a mixture of different
lengths. The question addressed here is whether the symmetric finite-volume
state selects a single length, and whether that length is the
magnetization selected by an infinitesimal field.

\subsection{Model and result}
Fix $S\in\{\tfrac12,1,\tfrac32,\ldots\}$. Let $S^1,S^2,S^3$ be the
irreducible spin-$S$ matrices, normalized by
\[
 [S^a,S^b]=i\sum_c\epsilon_{abc}S^c,
 \qquad \sum_{a=1}^3(S^a)^2=S(S+1)I.
\]
For even $L\ge4$, write $\Lambda_L=(\Z/L\Z)^3$ and $V=L^3$. On
$\bigotimes_{x\in\Lambda_L}\C^{2S+1}$, consider
\begin{equation}\label{eq:model}
 H_{S,L}=-\sum_{\{x,y\}\text{ nearest neighbors}}
              \sum_{a=1}^3 S_x^aS_y^a,
 \qquad M_L^a=\sum_{x\in\Lambda_L}S_x^a,
\end{equation}
where each unordered bond is counted once. We use the Gibbs expectation
\[
 \langle A\rangle_{\beta,L}
 =\frac{\Tr(Ae^{-\beta H_{S,L}})}{\Tr e^{-\beta H_{S,L}}}.
\]
The pressure and its right derivative at zero field are
\begin{align}
 p_{S,\beta}(h)
   &=\lim_{\substack{L\to\infty\\L\text{ even}}}
       \frac1{\beta V}\log\Tr
          e^{-\beta(H_{S,L}-hM_L^3)},\label{eq:pressure}\\
 m_{S,\beta}&=p_{S,\beta}'(0+).
       \label{eq:magnetization}
\end{align}
The pressure limit exists for this bounded finite-range interaction, and
convexity gives the right derivative. The volume limit in
\eqref{eq:pressure} precedes removal of the field.

\begin{theorem}[Spherical magnetization law]\label{thm:main}
For each fixed $S$ there is a finite $\beta_0(S)$ such that, for every
fixed finite $\beta\ge\beta_0(S)$, one has $m=m_{S,\beta}>0$ and
\begin{equation}\label{eq:spherical-law}
 \lim_{\substack{L\to\infty\\L\text{ even}}}
  \left\langle
     \exp\!\left(\frac1V\sum_{a=1}^3t_aM_L^a\right)
  \right\rangle_{\beta,L}
   =\frac1{4\pi}\int_{\mathbb S^2}e^{m t\cdot n}\dd\sigma(n)
   =\frac{\sinh(m|t|)}{m|t|}
 \qquad(t\in\R^3).
\end{equation}
Here $\sigma$ is surface-area measure on $\mathbb S^2$, and the quotient
at $t=0$ is interpreted as $1$. In addition,
\begin{equation}\label{eq:second-moment-law}
 \lim_{\substack{L\to\infty\\L\text{ even}}}
   \frac1{V^2}\sum_{x,y\in\Lambda_L}
          \langle S_x\cdot S_y\rangle_{\beta,L}=m^2.
\end{equation}
\end{theorem}

The statement uses self-adjoint linear combinations of the spin
components. It requires no joint finite-volume measurement of these
noncommuting operators. Both the spin and the positive temperature remain
fixed throughout the volume limit.

\subsection{Context and antecedents}\label{subsec:history}
The dimensional distinction in this problem is fundamental. Mermin and
Wagner~\cite{MerminWagner1966} excluded spontaneous ferro- or
antiferromagnetic order at positive temperature for finite-range
isotropic Heisenberg models in dimensions one and
two. In higher dimensions, infrared bounds proved continuous-symmetry
breaking for classical models~\cite{FSS1976} and ordering for quantum
antiferromagnets~\cite{DLS1978}. These quantum antiferromagnetic results
do not supply the isotropic ferromagnetic ordering input used here.
Another line of work establishes the low-temperature free-energy
asymptotic predicted by spin-wave theory~\cite{CGS2015}. That asymptotic
concerns a different limit from the field derivative or the zero-field
magnetization distribution at a fixed positive temperature.

The spherical distribution itself has clear antecedents. Ueltschi's
description of three-dimensional loop models~\cite[Section~8.1]{Ueltschi2017}
derives the transform $\sinh(hm)/(hm)$ for the spin-$\tfrac12$
ferromagnet conditional on a conjectured Poisson--Dirichlet law for
macroscopic loops. Bj\"ornberg, Fr\"ohlich, and
Ueltschi~\cite[Theorem~2.1]{BFU2020} prove the spherical transform
for the complete-graph quantum Heisenberg model at arbitrary fixed spin.
Our setting has nearest-neighbor interactions on the three-dimensional
lattice. We prove its magnetization law without assuming a distribution
for macroscopic loop lengths.

The existence of spontaneous magnetization and the identification of its
zero-field distribution are distinct questions. The companion manuscript
\cite{Ord} constructs magnetized equilibrium states for the isotropic
nearest-neighbor quantum Heisenberg ferromagnet in dimensions at least
three, for every fixed spin at sufficiently low positive temperatures.
The companion \cite{BL} establishes the low-temperature asymptotic
of the pressure-defined magnetization for nonnegative symmetric
finite-range interactions whose support generates $\Z^3$.
We use finite exchange and geometric estimates from these works. The present
argument adds spatial control on periodic tori and a pressure argument
excluding translation-ergodic equilibrium states with different
magnetization magnitudes.

The exchange representation builds on permutation and random-walk
representations of quantum spins~\cite{Powers1976,ConlonSolovej1991Representation},
T\'oth's random-stirring representation~\cite{Toth1993}, and its
higher-spin formulation through symmetric slots~\cite{Nachtergaele1994}.
The finite pin estimates used below draw on the stability and negative
dependence theory of Borcea, Br\"and\'en, and Liggett~\cite{BBL2009}.
Our zero-free matrix products use the contraction principle associated
with Asano's and Ruelle's developments of Lee--Yang theory
\cite{Asano1970,Ruelle1971,Ruelle1973}; the contraction needed here is
proved explicitly. Mean entropy, equilibrium states, and pressure
tangents are treated within the variational framework of Lanford and
Robinson~\cite{LanfordRobinson1968Entropy,LanfordRobinson1968Spin}.
We recall the required facts and prove the additional replica and
harmonic-selection statements.

\subsection{The two parts of the proof}\label{subsec:strategy}
Here is the obstruction that the proof must overcome. Let $J_L$ denote the
total-spin label in the rotation-invariant Gibbs state. Conditional on
$J_L$, any component of $M_L$ is uniform on
$-J_L,-J_L+1,\ldots,J_L$. A fixed positive field and the definition of $m$
already imply
\begin{equation}\label{eq:overview-upper}
 \PP(J_L/V>m+\epsilon)\longrightarrow0
 \qquad(\epsilon>0).
\end{equation}
Thus it remains to prove
\begin{equation}\label{eq:overview-lower}
 \liminf_{L\to\infty}\frac{\E[J_L(J_L+1)]}{V^2}\ge m^2.
\end{equation}
Together these bounds force $J_L/V\to m$ in probability and give
\eqref{eq:spherical-law}. Section~\ref{sec:global-spin} supplies this
last representation-theoretic argument in detail.

Two facts produce the lower bound. First, the squared magnetization
averaged on large local boxes is controlled by the squared magnetization
averaged on the whole torus, with an error vanishing as the local scale
grows. This prevents local magnetization from canceling through spatial
variation on larger scales. Second, every translation-ergodic equilibrium
state has magnetization vector of length $m$. Here an equilibrium state
maximizes mean entropy minus $\beta$ times the physical energy density
among translation-invariant
states; an ergodic state is an extreme point of that invariant state
space. Their variational and averaging properties are recalled in
Section~\ref{sec:equilibrium}. The second fact identifies the local
magnetization length to which the first fact applies.

Both facts use the same weak exponential estimate. To describe it, let
$\mathcal P$ be a rectangular partition whose cell sides are comparable
to a large scale $R\le L$. For a deterministic site function $f$ with
$|f_x|\le1$ and $\sum_{x\in C}f_x=0$ in each cell, put
$D_f=\sum_xf_xS_x^3$. We prove, for
$b_R\asymp R^{-5/2}$ and a fixed $\kappa>0$, that
\begin{equation}\label{eq:overview-exp}
 \log\langle e^{\theta_R D_f}\rangle_{\beta,L}
       \le C_\beta b_R V,
 \qquad \theta_R=\kappa b_R\log R.
\end{equation}
In particular the left side divided by $\theta_R V$ tends to zero
uniformly in these tests. In nested rectangular partitions, apply this
bound to signs of differences between a cell average and its parent
average. A union bound and a telescoping identity then control all spatial
scales. Applied to a product of two local zero-sum spin tests in
different replicas, it forces the expectation of that product to vanish
in every equilibrium state of the replicated system: finite-volume
restrictions of such a state have relative entropy per site tending to
zero with respect to the product Gibbs state.

Obtaining \eqref{eq:overview-exp} is the spatial part of the proof.
Sections~\ref{sec:exchange}--\ref{sec:spatial-comparison} use exchange
cycles and a field of strength $b_R$ in a perpendicular direction to bound
$R^{-5/2}V^{-1}\langle D_f^2\rangle_{b_R}$ by a negative power of $R$.
The field gives a killed one-particle kernel an invertible resolvent.
A sparse-pin estimate controls the holes left after exposing unpinned
cycles; a comparison on this random set then uses the zero sum of $f$ in
each cell. Section~\ref{sec:zero-field} turns the resulting variance
bound into \eqref{eq:overview-exp} through a zero-free complex strip.

The remaining argument identifies the magnitude. Negative squares of
large-box magnetization averages would distinguish states of different
lengths, but the corresponding perturbed Gibbs operators do not directly
fit the exchange zero-freeness argument. Section~\ref{sec:pressure}
realizes these perturbations as first-order averages of random products
of real matrices acting on the $2S$ spin-$\tfrac12$ slots representing
each site. Let $Z(z)$ be the normalized trace of such a product with
complex field $z$. For every realization, $Z(z)$ is nonzero on
$\Rea z>0$. Shared Gaussian variables couple its moments to replicated
pressures; the zero-sum cross terms just described disappear from their
pressure derivatives. Moment bounds then select harmonic limits of
normalized $\log|Z(z)|$. Analyticity of their restrictions to
the real axis excludes coexistence of different ergodic magnetization
lengths. More concretely, an ergodic state of length $a$, oriented along
the auxiliary field, contributes the affine function $xa-a^2$ to a
limiting pressure tangent. Large positive $x$
selects the branch $xm-m^2$; analytic continuation of that branch cannot
dominate a competing branch of smaller length at every positive $x$.
The spatial estimate, the replica argument, and the passage from
random zero-free functions to harmonic pressure tangents are the main
additional constructions of this paper.

\subsection{Conventions}\label{subsec:conventions}
We suppress $S$ in expectations and usually suppress $\beta,L$ when no
confusion is possible. Write $X,Y,Z$ for the site spin components
$1,2,3$, in the standard representation where $X,Z$ are real and $Y$ is
imaginary. Set $H_0=-\beta H_{S,L}$. Constants $c,C$ may depend on the fixed
spin and on fixed geometric choices; $C_\beta$ may also depend on the
fixed inverse temperature. We first choose a sufficiently large finite
$\beta$, and only then choose lower thresholds for $L$ and spatial scales.
All torus volume limits are through even $L$.

The finite exchange inputs are developed in Section~\ref{sec:exchange},
and the deterministic lattice estimates imported from the companions
are stated with their hypotheses where they are used. Positivity of
$m$ follows directly from our torus sparse-pin estimate in
Corollary~\ref{cor:positive-magnetization}; no sharp low-temperature
asymptotic is needed. We increase $\beta_0(S)$ as needed below.
